%!TEX TS-program = xelatex
%!TEX encoding = UTF-8 Unicode
%===============================================================================
%  MASTER CV  —  TECH / FAANG track
%  Target roles: ML Engineer, SWE-ML, Applied Scientist (ML), GenAI Engineer
%  READABILITY NOTES:
%   - Visible bullets = one line each, minimal bold (numbers + one marquee term).
%   - Deeper specifics live in "% INTERVIEW PREP" comments per entry.
%   - This is a *master* resume; trim/reorder per JD.
%===============================================================================
\documentclass[11pt, a4paper]{awesome-cv}
\geometry{left=1.4cm, top=1.35cm, right=1.4cm, bottom=1.35cm}
\colorlet{awesome}{darktext}
\usepackage{hyperref}

% ---- Tighter vertical spacing than stock, without cramping ----
\renewcommand{\acvSectionTopSkip}{2.2mm}
\renewcommand{\acvSectionContentTopSkip}{3.0mm}
\renewcommand{\acvHeaderAfterSocialSkip}{4.0mm}
\setlength{\parskip}{6pt}

\name{Ishaan}{Gupta}
\position{ML Engineer {\enspace}|{\enspace} Deep Learning Researcher}
\address{San Diego / Bay Area, CA}
\email{i3gupta@ucsd.edu}
\homepage{ishaanSD.github.io/home}
\linkedin{IgAI}
\github{IshaanSD}

\begin{document}
\makecvheader

%-------------------------------------------------------------------------------
\begin{cvparagraph}
% CS PhD researcher in deep learning applications, building ML systems where success depends on more than the model: noisy measurements, limited labels, domain constraints, robust pipelines, and reliable evaluation.
ML PhD researcher in deep learning applications, building robust AI pipelines for noisy, limited-label real-world data, with experience spanning biotechnology, surgical analytics, recommendation systems, and data engineering.
\end{cvparagraph}

%-------------------------------------------------------------------------------
\cvsection{Education}
\begin{cventries}
\cventry
{Ph.D. Computer Science — Deep Learning for Genomics \quad | \quad M.S. Computer Science, earned en route}
{University of California, San Diego}
{La Jolla, CA}
{Sep 2023 -- Mar 2028 (exp.)}
{
\begin{cvitems}
\item {M.S. awarded Jun 2025 {\enspace}|{\enspace} GPA 3.92 {\enspace}|{\enspace} Charles Lee Powell Research Fellowship {\enspace}|{\enspace} Advisors: Tiffany Amariuta, Pavel A. Pevzner}
\end{cvitems}
}

\cventry
{B.S. Computer Science — Bioinformatics Specialization}
{University of California, San Diego}
{La Jolla, CA}
{Aug 2019 -- Jun 2022}
{
\begin{cvitems}
\item {GPA 3.92 {\enspace}|{\enspace} Magna Cum Laude {\enspace}|{\enspace} IEEE Eta Kappa Nu {\enspace}|{\enspace} Minor in Biology}
\end{cvitems}
}
\end{cventries}

%-------------------------------------------------------------------------------
\cvsection{Experience}
\begin{cventries}


  % =============================== DIPLOFORMER =================================
  % INTERVIEW PREP (not shown on CV):
  %  - 186M-param FlashAttention transformer; progressive unfreezing
  %    (backbone frozen epoch 1, then adapted).
  %  - Data: 400 individuals; EITHER 8,100 loci @4,096bp OR 2,700 loci @49kbp.
  %  - HW: mixed precision on 2×V100 or single H100; <100-sample runs -> 2 jobs
  %    co-located on one H100.
  %  - Nextflow DAG: CPU nodes (ingest/preprocess) -> GPU (train) -> light GPU (infer).
  %  - Fault tolerance: preempted Slurm jobs resume into the SAME W&B run.
  %  - Baselines Elastic Net + XGBoost. Drift detection stratified by ancestry and by
  %    biological features of genes/peaks (generalized to "subpopulations" on CV).
  %  - Predicts personalized molecular phenotypes (gene expression, chromatin
  %    accessibility) that generalize across loci AND across people.
  %  - Venues: Keystone 2025 (poster); ASHG 2026.
  \cventry
    {Graduate Student Researcher — Deep Learning (Dr. Tiffany Amariuta)}
    {Halıcıoğlu Data Science Institute, UCSD}
    {San Diego, CA}
    {Sep 2024 -- Present}
    {
      \begin{cvitems}
        \item {Built \textit{DiploFormer}, a Nextflow/\textbf{PyTorch} framework to adapt foundation models to personal genomes for disease mechanisms discovery.}
        \item {Fine-tuned \textbf{186M-parameter} FlashAttention \textbf{transformer} + CNN-based model across people and sequence contexts on multi-modal data.}
        \item {Implemented delayed pairwise-loss computation via output accumulation, ran \textbf{WandB} / MLflow sweeps of adaptation strategies and ablations.}
        \item {Optimized resource-aware CPU/GPU workflows with parallel jobs, model checkpoint recovery across preprocessing, training, and inference.}
        \item {Benchmarked against Elastic Net/\textbf{XGBoost} baselines, monitored data drift, verified discovery of causal features through external validation.}
      \end{cvitems}
    }

  % =============================== INTUITIVE ===================================
  % INTERVIEW PREP (not shown on CV):
  %  - LOF (Local Outlier Factor) ~O(n^2), memory-heavy. Replaced with GMM (linear)
  %    or boxplot/IQR (O(n log n)). Fill [N]× after benchmarking on one batch.
  %  - Latency = BATCH/compute (not endpoint): wall-clock of outlier stage before/
  %    after; full analytics-refresh time; memory drop. Serving is FastAPI but
  %    team-owned — mention only if role centers on serving.
  %  - "Trend algorithms + significance testing vs. downstream outcomes."
  %  - Product (DO NOT reveal): which OR / prep time / surgeon is slow; scheduling
  %    utilization. "Robust diagnostics" = median/IQR/MAD summaries.
  \cventry
    {AI Data Scientist Intern}
    {Intuitive Surgical, Inc.}
    {Sunnyvale, CA}
    {Jun 2026 -- Sep 2026}
    {
      \begin{cvitems}
        \item {Engineered Python/SQL components across the full analytics pipeline, transforming action-recognition outputs from 100K+ surgical cases into validated event tables, benchmark metrics, and dashboard-ready \textbf{user-facing insights} for analyzing hospital operational efficiency.}
        \item {Improved \textbf{temporal analytics} robustness by evaluating \textbf{outlier-detection} strategies and implementing a custom GMM/IQR detector.}
        % \item {Built an LLM-assisted data-ingestion workflow to standardize inconsistent spreadsheet schemas from multi-site hospital observations, using rule-grounded prompts and human-in-the-loop validation to reconcile metric names and quantify longitudinal outcome trends.}
        \item {Built an LLM-assisted data-ingestion and validation workflow to standardize inconsistent hospital observation spreadsheets for user research}
        \item {Enabled detector robustness and longitudinal trend analyses that identified service-associated improvements in business outcomes.}
        % \item {Implemented validation and evaluation workflows measuring detector robustness, false-positives, and association with downstream KPIs.}
      \end{cvitems}
    }

  % =============================== ABTERRA ====================================
  % INTERVIEW PREP: 4× = improvement on UNSEEN, company-specific data after
  % fine-tuning autoregressive transformers for de novo peptide sequencing (mass-spec).
  \cventry
    {Machine Learning Engineer Intern}
    {Abterra Biosciences}
    {San Diego, CA}
    {Jul 2024 -- Sep 2024}
    {
      \begin{cvitems}
        \item {Benchmarked autoregressive transformer models for de novo peptide sequencing, improving held-out performance \textbf{4$\times$} on in-house data.}
        \item {Built reproducible PyTorch Lightning + W\&B training workflows for hyperparameter sweeps, model comparison, and failure-mode analysis.}
        \item {Achieved \textbf{95\% accuracy} (baseline: 60\%) in resolving ambiguous antibodies through peptide–spectrum feature engineering and FNN classifiers.}
        \item {Deployed trained deep learning models in the company's Java codebase using DL4J, enabling parallelized performance evaluation.}
      \end{cvitems}
    }
    
  % =============================== PEVZNER ====================================
  % INTERVIEW PREP:
  %  - Repeat-finding project: approximate + suffix-array pattern matching with an
  %    ITERATIVE de Bruijn graph simplification procedure; repeat-motif lengths
  %    10,000 to 1,000,000 characters. Beat SOTA repeat detection in the critical
  %    immune-system DNA loci (Primate T2T; Nature 2025). Suffix-array alignment
  %    (unialigner) optimized 4-10× in C++.
  %  - Signal-processing (SEPARATE project): cumulative sums (CUSUM) for nanopore
  %    blockade detection; FFT + clustering to infer directionality of the peptide
  %    sequence.
  %  - Papers: Nature (ape genomes / Primate T2T) + Bioinformatics (GenomeDecoder).
  %    Plotly genome-scale viz lives on the Bio CV.
  \cventry
    {Graduate Student Researcher — Algorithms (Dr. Pavel A. Pevzner)}
    {UCSD CSE Department}
    {San Diego, CA}
    {Jan 2022 -- Jul 2024}
    {
      \begin{cvitems}
        \item {Co-Designed GenomeDecoder (\textbf{C++}) for finding 10K-1M long inexact repeats in 100M long sequences by iteratively simplifying de Bruijn graphs}
        \item {Outperformed SOTA on segmental duplication detection in immune loci; Browser-friendly genome-scale visualization of the duplications \textit{Bioinformatics} (2025).}
        \item {Improved UniAligner (suffix-array-based alignment) runtime \textbf{4--10$\times$} and integrated with GenomeDecoder to discover gene deletions across primate genomes. \textit{Nature (2025)}}
        \item {Developed signal-processing methods using cumulative sums, FFT, and clustering for delineating bidirectionality in noisy Nanopore signals.}
      \end{cvitems}
    }
    
  % =============================== ILLUMINA ===================================
  % INTERVIEW PREP: Agile; LIMS = Lab Information Management System; genotyping
  % 10,000s samples/day; Java Spring + SQL + Angular; REST, AWS, RabbitMQ; containerized.
  \cventry
    {Software Engineer I — Systems Integration}
    {Illumina, Inc.}
    {San Diego, CA}
    {Jun 2022 -- Dec 2022}
    {
      \begin{cvitems}
        \item {Developed production-grade ETL infrastructure code for high-throughput lab automation workflows processing \textbf{10{,}000s of samples/day}.}
        \item {Built REST APIs and event-driven microservices using Java Spring/JDBC, SQL databases, RabbitMQ, AWS S3/SQS, and TypeScript.}
        \item {Deployed \textbf{Dockerized} services on \textbf{Kubernetes} with SQL stored procedures, JUnit tests, and CI/CD release workflows.}
      \end{cvitems}
    }
    % \cventry
    % {Drug Discovery Project Consultant}
    % {Model Medicines}
    % {San Diego, CA}
    % {Nov 2020 – Mar 2021}
    % {
    %   \begin{cvitems}
    %     \item {Engineered \textbf{GCP BigQuery}-based COVID-19 drug screening pipeline supporting \textbf{SQL} queries on ChEMBL data, and RDKit molecular fingerprints.}
    %     \item {Led project team for \textbf{NLP tool} to evaluate therapeutic potential by \textbf{text-mining} scientific journals for input drug keywords and their synonyms.}
    %   \end{cvitems}
    % }
\end{cventries}

%-------------------------------------------------------------------------------
\cvsection{Selected Projects}
\begin{cventries}
  % CHEMBOT — 3 MCP servers over openFDA (labels + FAERS) and PubMed; LangGraph
  % multi-route router; ChromaDB + SentenceTransformers + OpenAI; deployed on GCP.
    \cventry
    {LangGraph · MCP · RAG · ChromaDB · Google Cloud}
    {{ChemBot - Agentic RAG Assistant for Trusted OTC Drug Information}}
    {~}
    {(github.com/IshaanSD/ChemBot)}
    {
        \begin{cvitems}
            \item {Agentic system for over-the-counter drug information retrieval (eg: side-effects) grounded in trusted biomedical evidence (FAERS and PubMed).}
            \item {\textbf{LangGraph} routing across \textbf{MCP servers} integrated with RAG (ChromaDB), embeddings (SentenceTransformer), and LLMs (OpenAI/Gemini)}
            \item {Deployed on Google Cloud with automated evaluations for tool-calling, agent skills, citation faithfulness, and cross-LLM response quality.}
        \end{cvitems}
    }

  % GENERATIVE MODELS — latent diffusion (UNet, conditional chest X-ray), β-VAE
  % (interpretable latents), UNet-GAN (mammogram augmentation + tumor classification).
  \cventry
    {Diffusion Model · $\beta$-VAE · Context U-Net · GAN}
    {Generative Models for Medical Image Augmentation}
    {~}
    {(github.com/IshaanSD/DLCXR)}
    {
        \begin{cvitems}
            \item {Built a \textbf{diffusion model} training pipeline using $\beta$-VAE latent encoding and a Context U-Net denoiser, and demonstrated on chest X-ray images.}
            \item {Implemented disentangled $\beta$-VAE representations for CXR generation and a U-Net GAN baseline for mammography augmentation.}
        \end{cvitems}
    }


  % MITOEDIT — MCP server lets any agent run tool via NL; multi-format ingestion;
  % downstream queries e.g. "targets with no non-benign bystanders". Existing stack:
  % Python + FastAPI web UI + Docker + tests + alignment/liftover -> CSVs. bioRxiv preprint.
  \cventry
    {MCP · LLM Agents · FastAPI · Docker · Python}
    {MitoEdit — Lab-in-the-Loop Agent for Mitochondrial Research at St. Jude}
    {~}
    {(github.com/xz-stjude/MitoEdit)}
    {
        \begin{cvitems}
            \item {Built an \textbf{MCP server} and FastAPI layer that lets bench scientists run TalenWF through typed, agent-callable tools. \textit{(manuscript in preparation)}}
            \item {Dockerized a CI/CD-enabled Python workflow + UI for identifying mtDNA cancer variants editable by TALE and predicting bystander edits.}
        \end{cvitems}
    }
    % GRAPHAZON — recommendation / ranking. Learn-to-rank via LambdaMART (LightGBM);
    % NDCG/MRR; Amazon ESCI Shopping Queries; KG + Node2Vec + text embeddings.
    \cventry
    {LangChain · Knowledge Graphs · LightGBM / LambdaMART}
    {Graphazon - Natural Language Query-based Product Recommendation}
    {~}
    {(github.com/IshaanSD/Graphazon)}
    {
      \begin{cvitems}
        \item {Built a product-recommendation pipeline on the Amazon ESCI dataset, using an LLM to parse queries into typed intents.}
        \item {Combined knowledge-graph, Node2Vec, and text embeddings with \textbf{LambdaMART} learn-to-rank optimized for Mean Reciprocal Rank.}
      \end{cvitems}
    }
\end{cventries}

%-------------------------------------------------------------------------------
\cvsection{Technical Skills}
\begin{cvskills}
  \cvskill
    {Languages}
    {Python, Java, C++, SQL, Bash, R}
  \cvskill
    {ML / Deep Learning}
    {PyTorch, Lightning, Transformers, CNN, LSTM, PEFT / LoRA, FlashAttention, scikit-learn, XGBoost, LightGBM, TensorFlow}
  \cvskill
    {MLE / MLOps}
    {W\&B, MLflow, Hugging Face, Docker, FastAPI, REST APIs, CI/CD, pytest/JUnit, model evaluation, drift analysis}
  \cvskill
    {Data \& Infrastructure}
    {Spark, Kafka, Airflow, Pandas, Polars, SQL optimization, MongoDB, AWS, GCP, Slurm, Nextflow, Kubernetes}
  \cvskill
    {GenAI / Retrieval}
    {LangGraph, LangChain, MCP, RAG, tool-calling, ChromaDB, SentenceTransformers, agent evaluation, }
\end{cvskills}

%-------------------------------------------------------------------------------
\cvsection{Publications \& Conferences}
\begin{cventries}
\cventry
{~}{~}{~}{~}
{
\begin{cvitems}
% \item {\textbf{Gupta I.}, Amariuta T. \textit{DiploFormer: Deep Learning framework for capturing Non-Additive Regulatory Effects on Gene Expression and Chromatin Accessibility.} Presented at Keystone AI in Molecular Biology, 2025; \textbf{ASHG}, 2026. (github.com/IshaanSD/DiploFormer\_1000G\_atac)}
% \item {Zhang Z., \textbf{Gupta I.}, Pevzner P.A. \textit{GenomeDecoder: inferring segmental duplications in repetitive genomic regions.} \textbf{Bioinformatics}, 2025.}
% \item {Yoo D.A., et al. \textit{Complete sequencing of ape genomes.} \textbf{Nature}, 2025.} (ishaansd.github.io/iPlots)
\item {\textit{DiploFormer: Deep Learning framework for capturing Non-Additive Regulatory Effects on Gene Expression and Chromatin Accessibility.} Presented at Keystone AI in Molecular Biology, 2025; \textbf{ASHG}, 2026. (github.com/IshaanSD/DiploFormer\_1000G\_atac)}
\item {\textit{GenomeDecoder: inferring segmental duplications in repetitive genomic regions.} \textbf{Bioinformatics}, 2025.}
\item {\textit{Complete sequencing of ape genomes.} \textbf{Nature}, 2025.} (ishaansd.github.io/iPlots)
\end{cvitems}
}
\end{cventries}

\cvsection{Certifications}
\begin{cventries}
\cventry
{~}{~}{~}{~}
{
\begin{cvitems}
\item {\textbf{AI Agents:} Google AI Agents Capstone \quad | \quad \textbf{Data Engineering:} IBM SQL, ETL/Airflow/Kafka, Spark/Hadoop \quad | \quad \\ \textbf{Backend:} Advanced Java Spring Integration Testing}
\end{cvitems}
}
\end{cventries}

\cvsection{Teaching Experience}
\begin{cventries}
\cventry
{Teaching Assistant / Tutor}
{UC San Diego}
{San Diego, CA}
{Sep 2020 -- Jun 2022}
{
\begin{cvitems}
\item {Courses: \textit{CSE 181}, \textit{CSE 182}, \textit{BICD 140}, \textit{BILD 4}; taught algorithms, machine learning concepts, statistics for research, and biological data analysis.}
\item {Supported instruction on graph algorithms, dynamic programming, HMMs, clustering, sequence analysis, databases, and hypothesis testing.}
\end{cvitems}
}
\end{cventries}

% \cvsection{Teaching Experience}
% \begin{cventries}
% \cventry
% {Teaching Assistant / Tutor}
% {UC San Diego}
% {San Diego, CA}
% {Sep 2020 -- Jun 2022}
% {
% \begin{cvitems}
% \item {Taught data structures and algorithms, machine learning, biological databases, immunology, and lab methods across 4 courses, supporting technical instruction, grading, and student mentoring.}
% \item {Explained complex topics including graph algorithms, dynamic programming, HMMs, clustering, statistics, and scientific writing to undergraduate students.}
% \end{cvitems}
% }
% \end{cventries}


\end{document}
